% Appendix A: Hyperparameters, extra result tables, dataset stats, worked examples
\selectlanguage{greek}


\chapter{Συμπληρωματικά Αποτελέσματα και Στατιστικά Συνόλου Δεδομένων}
\label{ch:a}


\section{Διαμορφώσεις Υπερπαραμέτρων}
\label{sec:hyperparams}

Οι πίνακες~\ref{tab:hyperparam-classifiers} και~\ref{tab:hyperparam-features} καταγράφουν
τις υπερπαραμέτρους όλων των ταξινομητών και των pipeline εξαγωγής χαρακτηριστικών που
χρησιμοποιούνται σε αυτή την εργασία.

Οι ταξινομητές Λογιστικής Παλινδρόμησης για το Επίπεδο~1, το Επίπεδο~2 (ενσωμάτωση) και
τον αυτόνομο \texttt{\en{TfidfClassifier}} μοιράζονται ακριβώς τις ίδιες προεπιλεγμένες τιμές,
επιλεγμένες για ευρωστία και αναπαραγωγιμότητα χωρίς \en{fine-tuning}.
Ο \texttt{\en{PosClassifier}} χρησιμοποιεί \en{Random Forest} με \texttt{\en{class\_weight=`balanced'}}
επειδή η κατανομή κλάσεων στα χαρακτηριστικά \en{POS} είναι άνιση.
Ο \texttt{\en{SvmClassifier}} εκπαιδεύεται σε στρωματοποιημένο δείγμα \SvmSampleSize{} ζευγών
λόγω του τετραγωνικού έως κυβικού κόστους του \en{SVC} με πυρήνα ακτινικής βάσης.

\begin{table}[h!]
  \centering
  \small
  \begin{tabular}{llll}
    \toprule
    \textbf{Μοντέλο} & \textbf{Αλγόριθμος} & \textbf{Παράμετρος} & \textbf{Τιμή} \\
    \midrule
    \texttt{\en{StructuralClassifier}} & \en{Logistic Regression} & $C$ & \StructuralC{} \\
                                       &                           & \texttt{\en{max\_iter}} & \StructuralMaxIter{} \\
                                       &                           & \texttt{\en{random\_state}} & \Seed{} \\
    \midrule
    \texttt{\en{TfidfClassifier}}      & \en{Logistic Regression} & $C$ & \TfidfClfC{} \\
                                       &                           & \texttt{\en{max\_iter}} & \TfidfClfMaxIter{} \\
                                       &                           & \texttt{\en{random\_state}} & \Seed{} \\
    \midrule
    \texttt{\en{PosClassifier}}        & \en{Random Forest}       & \texttt{\en{n\_estimators}} & \PosNEstimators{} \\
                                       &                           & \texttt{\en{class\_weight}} & \texttt{\en{`\PosClassWeight{}'}} \\
                                       &                           & \texttt{\en{random\_state}} & \Seed{} \\
    \midrule
    \texttt{\en{EmbeddingClassifier}}  & \en{Logistic Regression} & $C$ & \EmbeddingC{} \\
                                       &                           & \texttt{\en{max\_iter}} & \EmbeddingMaxIter{} \\
                                       &                           & \texttt{\en{random\_state}} & \Seed{} \\
    \midrule
    \texttt{\en{SvmClassifier}}        & \en{SVC (RBF kernel)}    & $C$ & \SvmC{} \\
                                       &                           & $\gamma$ & \texttt{\en{`\SvmGamma{}'}} \\
                                       &                           & \texttt{\en{subsample\_size}} & \SvmSampleSize{} \\
                                       &                           & \texttt{\en{random\_state}} & \Seed{} \\
    \midrule
    \en{CascadeLP}                     & ---                       & $\tau_1$ & \TauOneDefault{} \\
                                       &                           & $\tau_2$ & \TauTwoDefault{} \\
    \bottomrule
  \end{tabular}
  \caption{Υπερπαράμετροι ταξινομητών.}
  \label{tab:hyperparam-classifiers}
\end{table}

\begin{table}[h!]
  \centering
  \small
  \begin{tabular}{lll}
    \toprule
    \textbf{Pipeline} & \textbf{Παράμετρος} & \textbf{Τιμή} \\
    \midrule
    \en{TF-IDF} & \texttt{\en{max\_features}} & \TfidfMaxFeatures{} \\
                & \texttt{\en{sublinear\_tf}} & \texttt{\en{\TfidfSublinearTf{}}} \\
                & \texttt{\en{min\_df}} & \TfidfMinDf{} \\
    \midrule
    \en{Sentence-Transformer} & μοντέλο & \texttt{\en{\EmbeddingModelName{}}} \\
    \midrule
    \en{Node2Vec} & \texttt{\en{dimensions}} & \NTwoVDimensions{} \\
                  & \texttt{\en{walk\_length}} & \NTwoVWalkLength{} \\
                  & \texttt{\en{num\_walks}} & \NTwoVNumWalks{} \\
                  & \texttt{\en{window}} & \NTwoVWindow{} \\
                  & $p$ & \NTwoVP{} (αμερόληπτο \en{DeepWalk}) \\
                  & $q$ & \NTwoVQ{} (αμερόληπτο \en{DeepWalk}) \\
    \bottomrule
  \end{tabular}
  \caption{Υπερπαράμετροι pipeline εξαγωγής χαρακτηριστικών.}
  \label{tab:hyperparam-features}
\end{table}


\section{Πρόσθετοι Πίνακες Αποτελεσμάτων}
\label{sec:extra-results}

Ο πίνακας~\ref{tab:full-model-comparison} παρουσιάζει τα αποτελέσματα επικύρωσης του
\en{CascadeLP} και του μοντέλου αναφοράς \en{TF-IDF + SVM} (\S\ref{sec:svm-results}) στις
τέσσερις μετρικές που ορίζει το πρωτόκολλο αξιολόγησης (\S\ref{subsec:evaluation-metrics}).
Η σύγκριση τοποθετείται σε ενιαίο πίνακα για ευκολία αναφοράς.

\begin{table}[h!]
  \centering
  \small
  \begin{tabular}{lrrrr}
    \toprule
    \textbf{Μοντέλο} & \textbf{\en{Macro F1}} & \textbf{\en{AUC-ROC}} & \textbf{F1 ($\mathrm{CN}=0$)} & \textbf{Καθυστ.\ (ms/ζ.)} \\
    \midrule
    \en{CascadeLP} (τελικό) & \CascadeValFone{} & \CascadeValAuc{} & \CascadeValFone{} & \ThroughputLatency{} \\
    \en{TF-IDF + SVM}       & \SvmFone{}        & \SvmAuc{}         & \SvmCsFone{}      & \SvmLatency{}      \\
    \bottomrule
  \end{tabular}
  \caption{Σύγκριση \en{CascadeLP} έναντι \en{SVM} σε όλες τις μετρικές αξιολόγησης,
  σύνολο επικύρωσης (n=\ValSplitSize{}).
  Το F1 μηδενικού \en{CN} του \en{CascadeLP} ταυτίζεται αριθμητικά με το συνολικό
  \en{Macro F1}, επειδή σχεδόν όλα τα ζεύγη επικύρωσης ανήκουν σε αυτό το
  διαγνωστικό υποσύνολο. Η μάσκα δεν ταυτίζεται με το λειτουργικό
  \en{cold-start} της δρομολόγησης (\S\ref{sec:cold-start-results}).}
  \label{tab:full-model-comparison}
\end{table}

Ο πίνακας~\ref{tab:baseline-tier-summary} συνοψίζει την απόδοση των ταξινομητών ανά επίπεδο
ως αυτόνομα μοντέλα αναφοράς, όπως προκύπτει από την ανάλυση ανά επίπεδο στον
πίνακα~\ref{tab:cascade-tier-results} (Κεφάλαιο~\ref{ch:experiments}).
Σε αντίθεση με τη
σύγκριση του πίνακα~\ref{tab:full-model-comparison}, η οποία αφορά ανεξάρτητη εκπαίδευση
κάθε μοντέλου, εδώ αναφερόμαστε στην απόδοση του κάθε επιπέδου \textit{ως μέρος της
αλυσίδας}: δηλαδή στο υποσύνολο ζευγών που ο εκάστοτε ταξινομητής επέλυσε
ουσιαστικά.

\begin{table}[h!]
  \centering
  \small
  \begin{tabular}{llrrrr}
    \toprule
    \textbf{Επίπεδο} & \textbf{Ταξινομητής} & \textbf{n} & \textbf{Ποσοστό} & \textbf{Ακρίβεια} & \textbf{\en{Macro F1}} \\
    \midrule
    Επίπεδο~0 & αυτοβρόχοι        & 0              & 0{,}00\%         & ---            & ---            \\
    Επίπεδο~1 & \en{StructuralClassifier} & \TierOneCount{} & \TierOnePct\% & \TierOneAcc{} & \TierOneFone{} \\
    Επίπεδο~2 & \en{PosClassifier}        & \TierTwoCount{} & \TierTwoPct\% & \TierTwoAcc{} & \TierTwoFone{} \\
    Επίπεδο~3 & \en{EmbeddingClassifier}  & \TierThreeCount{} & \TierThreePct\% & \TierThreeAcc{} & \TierThreeFone{} \\
    \midrule
    \textbf{Σύνολο} & \en{CascadeLP} & \ValSplitSize{} & 100\% & --- & \CascadeValFone{} \\
    \bottomrule
  \end{tabular}
  \caption{Απόδοση ανά επίπεδο στο σύνολο επικύρωσης. Το χαμηλό \en{Macro F1} του
  Επιπέδου~1 οφείλεται στην ακραία ανισορροπία κλάσεων εντός του υποσυνόλου υψηλής
  αξιοπιστίας (\S\ref{sec:cascade-results}).}
  \label{tab:baseline-tier-summary}
\end{table}


\section{Στατιστικά Συνόλου Δεδομένων}
\label{sec:dataset-stats}

Ο πίνακας~\ref{tab:dataset-stats} συγκεντρώνει τα βασικά στατιστικά του \en{Wikipedia DSAA
2023}, όπως προκύπτουν από τους ελέγχους διαρροής και χαρακτηρισμού διαχωρισιμότητας
(Κεφάλαιο~\ref{ch:experiments}).

\begin{table}[h!]
  \centering
  \small
  \begin{tabular}{lr}
    \toprule
    \textbf{Μέγεθος} & \textbf{Τιμή} \\
    \midrule
    Ζεύγη εκπαίδευσης (\texttt{\en{train.csv}}) & \TrainPairs{} \\
    Ζεύγη ελέγχου (\texttt{\en{test.csv}})      & \TestPairs{}  \\
    Σύνολο άρθρων \en{Wikipedia} (\texttt{\en{nodes.tsv}}) & \GraphNodesTotal{} \\
    Κόμβοι στο δομικό γράφημα & \GraphNodes{} \\
    Μέσος βαθμός κόμβου & \GraphMeanDegree{} \\
    \midrule
    Κατανομή κλάσεων εκπαίδευσης (θετικά / αρνητικά) & \TrainPosPct\% / \TrainNegPct\% \\
    \midrule
    Αυτοβρόχοι εκπαίδευσης (θετικοί / αρνητικοί) & \SelfLoopsTrainPos{} / \SelfLoopsTrainNeg{} \\
    Αυτοβρόχοι ελέγχου & \SelfLoopsTest{} \\
    \midrule
    Ακριβής επικάλυψη εκπαίδευσης--ελέγχου & \LeakageExact{} ζεύγη \\
    Αντεστραμμένη επικάλυψη & \LeakageReversed{} ζεύγη \\
    Ενδο-συνόλου διπλοεγγραφή (μη κατευθυνόμενα ζεύγη) & \IntraTrainDuplicates{} \\
    \midrule
    Ζεύγη εκπαίδευσης χωρίς αυτοβρόχους & \TrainPairsNoSelf{} \\
    Σύνολο εκπαίδευσης (80\%) & \TrainSplitSize{} \\
    Σύνολο επικύρωσης (20\%) & \ValSplitSize{} \\
    \midrule
    Όριο $\theta_{\mathrm{CN}}$ (1\% FPR) & $> \CnThreshold{}$ \\
    Όριο $\theta_{\mathrm{TF\text{-}IDF}}$ (1\% FPR) & $> \TfidfThreshold{}$ \\
    \bottomrule
  \end{tabular}
  \caption{Στατιστικά συνόλου δεδομένων \en{Wikipedia DSAA 2023}. Τα όρια
  $\theta_{\mathrm{CN}}$ και $\theta_{\mathrm{TF\text{-}IDF}}$ υπολογίζονται με
  \texttt{\en{pick\_thresholds}} (\S\ref{sec:forensics-protocol}).}
  \label{tab:dataset-stats}
\end{table}

Αξίζει να σημειωθεί η χαμηλή τιμή του μέσου βαθμού κόμβου (\GraphMeanDegree{}).
Σε ένα τόσο αραιό γράφημα, η συντριπτική πλειονότητα των ζευγών στερείται κοινών γειτόνων,
γεγονός που εξηγεί γιατί το \ColdStartPct\% των ζευγών επικύρωσης ανήκει στο υποσύνολο μηδενικού \en{CN}
(\S\ref{sec:cold-start-results}) και γιατί τα ευρετικά χαρακτηριστικά δομής αποτυγχάνουν
στην πλειονότητα των ζευγών χωρίς \en{Node2Vec}.


\section{Σύνολο Ζευγών Ανάλυσης}
\label{sec:worked-examples}

Ο Πίνακας~\ref{tab:worked-examples} καταγράφει, ανά γραμμή, όλα τα ζεύγη που
συζητούνται ρητά στην ποιοτική ανάλυση των Επιπέδων~0--3 (\S\ref{subsec:tier-examples},
\S\ref{sec:tier3-hard-residual}): το αναγνωριστικό \texttt{\en{id}} τους στο
\texttt{\en{train.csv}}, ώστε κάθε ισχυρισμός να ανιχνεύεται σε συγκεκριμένη γραμμή
του \texttt{\en{cascade\_val\_tiers.csv}} (ή του \texttt{\en{train.csv}} για τους
αυτοβρόχους του Επιπέδου~0, οι οποίοι εξαιρούνται από την εκπαίδευση των
Επιπέδων~1--3 και άρα δεν εμφανίζονται στο πρώτο αρχείο), τα δύο άρθρα, το
επίπεδο που έδωσε την τελική πρόβλεψη, τον χαρακτηρισμό δυσκολίας
(\S\ref{sec:forensics-protocol}), την πραγματική και την προβλεπόμενη ετικέτα, και
το αν η πρόβλεψη ήταν ορθή.

\begin{table}[h!]
  \centering
  \footnotesize
  \resizebox{\textwidth}{!}{%
  \begin{tabular}{r p{3.6cm} p{3.6cm} c l c c c}
    \toprule
    \textbf{\en{id}} & \textbf{Άρθρο 1} & \textbf{Άρθρο 2} & \textbf{Επ.} & \textbf{Δυσκολία} & \textbf{Αλ.} & \textbf{Πρ.} & \textbf{Ορθή} \\
    \midrule
    557    & \en{Young Lake}            & \en{Young Lake} (αυτοσύνδεσμος)      & 0 & \en{self\_loop}       & 1 & 1 & Ναι \\
    289062 & \en{Hartbeespoort Aerial Cableway} & (αυτοσύνδεσμος)               & 0 & \en{self\_loop}       & 0 & 1 & Όχι \\
    \midrule
    769405 & \en{Molanna}               & \en{Molanna ulmerina}                & 1 & \en{high\_cn}         & 1 & 1 & Ναι \\
    310976 & \en{The Golden Mile}       & \en{Angling}                         & 1 & \en{high\_textsim}    & 1 & 1 & Ναι \\
    1021664& \en{Leaf River} (αποσαφ.)  & \en{Wapello} (αποσαφ.)               & 1 & \en{high\_textsim}    & 0 & 1 & Όχι \\
    479304 & \en{Viktoriya Mikhnovich}  & \en{Belarus}                         & 1 & \en{hard}             & 1 & 1 & Ναι \\
    947382 & \en{Plethodontohyla notosticta} & \en{The Phi Beta Kappa Society} & 1 & \en{hard}             & 0 & 1 & Όχι \\
    \midrule
    576964 & \en{N.\,G.\,Paulet, 18th Marquess of Winchester} & \en{Marquessate of Winchester} & 2 & \en{high\_textsim} & 1 & 1 & Ναι \\
    7608   & \en{Hood River} (αποσαφ.)  & \en{Hood River} (ποταμός)            & 2 & \en{high\_textsim}    & 1 & 0 & Όχι \\
    972868 & \en{No.\,3 Komarapalayam}  & \en{Mallur} (σταθμός)                & 2 & \en{hard}             & 1 & 1 & Ναι \\
    907028 & \en{Dutch Uncle} (ιδίωμα)  & \en{Dutch Uncle} (θεατρικό έργο)     & 2 & \en{hard}             & 1 & 0 & Όχι \\
    \midrule
    830738 & \en{FASTRAC 1}             & \en{JCB Fastrac}                     & 3 & \en{hard}             & 1 & 0 & Όχι \\
    505706 & \en{Abacetus laevigatus}   & (άσχετο είδος \en{Abacetus})         & 3 & \en{hard}             & 0 & 1 & Όχι \\
    841732 & \en{poor me} (αποσαφ.)     & \en{Victim mentality}                & 3 & \en{hard}             & 1 & 0 & Όχι \\
    735703 & \en{Abacetus laevigatus}   & \en{Tidal downsizing}                & 3 & \en{hard}             & 0 & 0 & Ναι \\
    710094 & \en{Niigata Institute of Technology} & \en{Kashiwazaki} (αποσαφ.) & 3 & \en{hard}             & 1 & 1 & Ναι \\
    37745  & \en{Stormrider}            & \en{David Gemmell}                   & 3 & \en{hard}             & 1 & 0 & Όχι \\
    429190 & \en{Stephen Breyer}        & \en{Individuals with Disabilities Education Act} & 3 & \en{hard} & 1 & 0 & Όχι \\
    \bottomrule
  \end{tabular}%
  }
  \caption{Πλήρες σύνολο ζευγών που αναλύονται ποιοτικά στο Κεφάλαιο~\ref{ch:analysis}
  (\S\ref{subsec:tier-examples}, \S\ref{sec:tier3-hard-residual}). Στήλες: επίπεδο
  \en{CascadeLP} που έδωσε την τελική πρόβλεψη (\en{Επ.}), χαρακτηρισμός
  δυσκολίας, πραγματική ετικέτα (\en{Αλ.}), προβλεπόμενη ετικέτα (\en{Πρ.}).}
  \label{tab:worked-examples}
\end{table}
